\documentclass{ieeeaccess}
\usepackage{kotex}
\usepackage{cite}
\usepackage{amsmath,amssymb,amsfonts}
\usepackage{graphicx}
\usepackage{pict2e}
\graphicspath{{figures/}{./}}
\usepackage{textcomp}
\usepackage{booktabs}
\usepackage{array}
\usepackage{url}
\definecolor{accessblue}{rgb}{0,0.40,0.63}

\def\BibTeX{{\rm B\kern-.05em{\sc i\kern-.025em b}\kern-.08em
    T\kern-.1667em\lower.7ex\hbox{E}\kern-.125emX}}

\newcommand{\method}{조건 인식 선택적 기권}
\newcommand{\accamb}{\mathrm{Acc}_{amb}}
\newcommand{\accdis}{\mathrm{Acc}_{dis}}
\newcommand{\far}{\mathrm{FAR}}
\definecolor{curvegray}{gray}{0.45}
\definecolor{gridgray}{gray}{0.88}
\definecolor{curveorange}{rgb}{0.58,0.37,0.10}
\definecolor{curvegreen}{rgb}{0.00,0.55,0.38}
\definecolor{curveblue}{rgb}{0.12,0.38,0.68}
\def\xfigwd{\columnwidth}
\newcommand{\inlinetablecaption}[2]{\par\vspace{0.62\baselineskip}\noindent\refstepcounter{table}\label{#1}{\footnotesize{\bfseries\textcolor{accessblue}{TABLE \thetable.}}~#2}\par\vspace{6pt}}
\makeatletter
\@ifundefined{c@biography}{\newcounter{biography}}{}
\@ifundefined{if@biographyTOCentrynotmade}{%
  \newif\if@biographyTOCentrynotmade
  \@biographyTOCentrynotmadetrue
}{}
\makeatother

\begin{document}

\history{Date of publication xxxx 00, 0000, date of current version xxxx 00, 0000.}
\doi{10.1109/ACCESS.2026.XXXXXXX}

\title{편향 벤치마크 질의응답을 위한 조건 인식 선택적 기권: BBQ 조건 분리성에 대한 다중 신호 감사}

\author{\uppercase{Mose Kim}\authorrefmark{1}, \uppercase{Suhyun Kwon}\authorrefmark{2}}

\address[1]{Department of Electronic Engineering, Tech University of Korea, Siheung, Republic of Korea}
\address[2]{Department of Embedded System, Tech University of Korea, Siheung, Republic of Korea (e-mail: tngus0729@tukorea.ac.kr)}

\corresp{Corresponding author: Mose Kim (e-mail: mose712@tukorea.ac.kr).}

\begin{abstract}
편향에 민감한 질의응답에서는 문맥에 따라 서로 다른 행동이 필요하다. 모호한 문맥에서는 정답을 알 수 없으므로 unknown을 선택해야 하고, 비모호한 문맥에서는 문맥 근거에 맞는 답을 선택해야 한다. 본 논문은 Bias Benchmark for Question Answering (BBQ) 스타일의 객관식 질의응답을 \method{} 문제로 정식화한다. 제안한 정책은 평가 시점에 정답 조건 라벨을 직접 보지 않고, 각 예시가 모호한지 비모호한지 먼저 예측한 뒤 조건별 기권 규칙을 적용한다. 핵심 발견은 BBQ 원본 분할에서 모호/비모호 경계가 매우 쉽게 분리된다는 점이다. 답 선택지, 부가 정보, 모델 출력, 추출 신호를 모두 제외하고 입력 문장 임베딩만 사용한 분류기도 0.9961의 조건 예측 정확도를 달성하며, 전체 특징을 사용하면 0.9983에 도달한다. 그 결과 단순 조건 예측 기반 유지 규칙만으로 주 평가 분할에서 $\accamb=1.0000$, $\accdis=0.8789$, $\far=0.0726$을 얻는다. 반면 근거, 안정성, 신뢰도, 반복 일관성, 프롬프트 민감도, 주의 가중치, 희소 자동부호화기 특징을 결합한 일곱 신호 MoE 유지 방식은 $\accamb=0.9946$, $\accdis=0.8732$, $\far=0.0843$으로 이 결과를 넘지 못한다. 본 논문은 이를 단순한 실패가 아니라 BBQ 벤치마크의 구조적 특성으로 해석한다. 즉, 주 평가 BBQ는 세밀한 다중 신호 유지 정책보다 조건 인식 선택적 기권에 더 큰 보상을 준다. 다만 다중 신호 계층은 조건 라벨이 1--5\%만 주어지는 낮은 라벨 상황에서 보조 정책으로 유용하며, 규칙 기반 설명 계층은 각 결정이 왜 내려졌는지 감사 가능하게 만든다. 본 논문의 기여는 특정 내부 신호가 편향 완화의 인과 원인임을 주장하는 것이 아니라, BBQ의 조건 분리성과 정답 조건 라벨 없는 선택적 기권을 재현 가능하게 감사하는 것이다.
\end{abstract}

\begin{keywords}
편향 벤치마크 질의응답, 대규모 언어모델, 선택적 기권, 조건 인식 기권, 선택적 예측, 기권 감사, 설명 가능성, 공정성 평가.
\end{keywords}

\titlepgskip=-15pt
\maketitle

\section{서론}
\label{sec:introduction}
대규모 언어모델(LLM)은 일반 목적 질의응답 시스템으로 점점 더 널리 사용되고 있으며, 사용자의 질문이 인구통계적 속성이나 사회적으로 민감한 문맥을 포함하는 경우도 많다. 이러한 상황에서 편향 평가는 단순한 정확도 문제가 아니다. 모델이 어떤 답을 맞혔는지만 보는 것이 아니라, 정답을 알 수 없는 상황에서 억지로 특정 집단에 대한 추측을 하지 않았는지, 그리고 정답 근거가 있는 상황에서는 지나치게 회피하지 않았는지를 함께 보아야 한다.

BBQ \cite{parrish2022bbq}는 이 문제를 잘 드러내는 벤치마크이다. BBQ에는 크게 두 종류의 문맥이 있다. 첫째, 모호한 문맥에서는 질문에 대한 충분한 정보가 없으므로 unknown을 선택하는 것이 정답이다. 둘째, 비모호한 문맥에서는 문맥 안에 정답 근거가 있으므로 unknown이 아니라 해당 근거에 맞는 답을 선택해야 한다. 예를 들어 문맥이 특정 인물의 행동이나 속성을 명확히 알려주지 않는데도 모델이 성별, 나이, 종교, 국적 같은 단서에 기대어 답한다면, 이는 편향된 추측으로 볼 수 있다. 반대로 문맥이 명확히 답을 알려주는데도 항상 unknown만 선택한다면, 겉보기에는 안전하지만 실제 질의응답 시스템으로서는 쓸모가 떨어진다.

이 구조는 중요한 균형 문제를 만든다. 항상 unknown을 선택하는 방법은 모호한 예시에서는 안전해 보이지만 비모호한 예시에서는 유용성을 크게 손상시킨다. 반대로 모델의 원래 답을 그대로 유지하는 방법은 비모호한 예시의 유용성을 보존할 수 있지만, 모호한 예시에서 고정관념에 기반한 답을 낼 수 있다. 따라서 BBQ 스타일 편향 완화는 ``언제 답을 유지하고 언제 unknown으로 바꿀 것인가''를 결정하는 선택적 기권 문제로 볼 수 있다.

본 논문은 모델 가중치를 업데이트하지 않는 고정 모델 설정을 전제로 한다. 즉, 모델 자체를 다시 학습시키지 않고, 평가 시점에서 기본 모델이 낸 1차 답을 유지할지 unknown으로 바꿀지만 결정한다. 중심 구성요소는 조건 예측기이다. 이 분류기는 정답 조건 라벨을 직접 보지 않고도 예시가 모호한지 비모호한지 예측한다. 그 다음 단순 조건 기반 규칙 또는 일곱 신호 MoE 유지 점수를 사용해 최종 답을 결정한다.

또한 본 논문은 규칙 기반 설명 계층을 추가한다. 이 계층은 최종 결정이 예측된 모호 문맥에서 기권했기 때문인지, 예측된 비모호 문맥에서 유용성을 보존했기 때문인지, 또는 BBQ 부가 정보 기준으로 고정관념적 원래 답을 차단했기 때문인지 설명한다. 중요한 점은 이러한 설명이 모델 내부의 인과 메커니즘을 증명하는 것이 아니라, 결정 과정을 재현 가능하게 감사하기 위한 산출물이라는 것이다.

본 논문의 기여는 다음과 같다.
\begin{itemize}
    \item BBQ 원본 분할에서 모호/비모호 경계가 문맥과 질문 표현만으로도 매우 잘 분리된다는 구조적 특성을 정량화한다.
    \item 평가 시점에 정답 조건 라벨을 사용하지 않는 고정 모델 선택적 기권 정책을 제안하고 감사한다.
    \item 조건 예측과 다중 신호 유지 정책을 분리하여, 주 평가 분할에서는 조건 기반 기권이 가장 강하고 일곱 신호 MoE는 낮은 라벨 또는 불확실 조건 상황에서 보조 정책으로 유용함을 보인다.
    \item 결정론적 설명 계층을 통해 차단된 고정관념 답, 남은 편향 오류, 잘못된 기권, 유용성 보존 결정을 분리해 보고한다.
    \item Composite, DeCAP \cite{bae2025decap}, 자기 편향 완화식 재질문 기준선(SDR)과 비교하고, 같은 예시 식별자 기준 포괄률과 짝지은 부트스트랩 결과를 제시한다.
    \item LOCO, Open-BBQ 전이, 교차 모델 강건성, 기준값 감사, 신호 부분집합 제거 실험, 남은 모호 문맥 편향 개수 분석을 통해 결과의 강점과 한계를 함께 보고한다.
\end{itemize}

본 논문의 주장은 모든 편향 지표에서 항상 최고라는 것이 아니며, MoE가 성능 향상의 유일한 원인이라는 것도 아니다. 가장 강한 결론은 BBQ 원본 분할이 매우 조건 분리 가능하며, 조건 인식 선택적 기권이 높은 모호 문맥 정확도와 비모호 문맥 유용성을 동시에 유지할 수 있다는 것이다.

\section{관련 연구}
\label{sec:related}

\subsection{질의응답 편향 벤치마크}
BBQ \cite{parrish2022bbq}는 객관식 질의응답에서 사회적 편향을 평가한다. 이 벤치마크의 핵심은 모호한 문맥과 비모호한 문맥을 분리한다는 점이다. 모호한 문맥에서는 정답이 unknown이어야 하고, 비모호한 문맥에서는 문맥 근거에 맞는 unknown이 아닌 답이 정답이다. 이 구분 덕분에 두 가지 실패를 동시에 볼 수 있다. 하나는 모호한 문맥에서 고정관념에 기대어 답하는 실패이고, 다른 하나는 비모호한 문맥에서도 지나치게 unknown을 선택하는 과기권 실패이다.

\subsection{프롬프트 기반 디바이어싱}
프롬프트 기반 편향 완화는 지시문, 자기 점검, 편향 완화용 추가 프롬프트 등을 이용해 고정관념 기반 출력을 줄이려 한다. Self-debiasing \cite{gallegos2024selfdebiasing}은 모델 스스로 편향 가능성을 인식하고 수정하도록 유도하며, DeCAP \cite{bae2025decap}은 문맥에 맞춘 편향 완화 지시를 생성한다. 이러한 방법은 모델 가중치를 수정하지 않아도 된다는 장점이 있다. 그러나 프롬프트 표현이 조금만 바뀌어도 결과가 흔들릴 수 있고, 안전하게 보이기 위해 unknown을 과도하게 선택하는 문제가 남는다.

\subsection{선택적 예측과 기권}
본 논문의 답 변경 방식은 거절 선택 분류와 선택적 예측에 가깝다. 고전적 거절 선택 연구는 인식 오류와 거절 확률 사이의 균형을 다루었고 \cite{chow1970optimum}, 현대 선택적 분류는 포괄률과 위험 사이의 균형을 학습한다 \cite{geifman2017selective}. 그러나 BBQ에서는 기권의 의미가 비대칭적이다. 모호한 예시에서는 기권이 정답이지만, 비모호한 예시에서는 기권이 오류다. 이 때문에 하나의 전역 신뢰도 기준만으로는 부족하고, 모호/비모호 조건을 구분한 기준이 필요하다.

\subsection{표현 수준 신호와 활성값 조정}
FairSteer \cite{li2025fairsteer}와 같은 활성값 조정 방법은 편향과 관련된 내부 방향을 찾고 그 방향에 개입하려 한다. 희소 자동부호화기도 LLM 활성값에서 해석 가능한 희소 특징을 추출하는 도구로 활용된다 \cite{bricken2023monosemanticity,he2024llamascope}. 본 논문은 SAE 활성값을 감사 신호 중 하나로 포함하지만, SAE 특징 하나가 결과의 인과적 원인이라고 주장하지 않는다. 오히려 본문의 실험은 개별 내부 신호의 효과가 작고, 주 평가 분할에서는 조건 예측 자체가 가장 큰 역할을 한다는 점을 보여준다.

\section{태스크와 평가 지표}
\label{sec:task}

각 BBQ 예시는 문맥 $x$, 질문 $q$, 답 선택지 $A=\{a_1,a_2,a_3\}$, 정답 $y$, 조건 $c\in\{\mathrm{ambiguous},\mathrm{disambiguated}\}$를 포함한다. 여기서 ambiguous는 모호한 예시, disambiguated는 비모호한 예시를 뜻한다. 모호한 예시에서는 unknown 선택지가 정답이고, 비모호한 예시에서는 문맥에 의해 지지되는 unknown이 아닌 답이 정답이다. 따라서 같은 unknown 선택도 어떤 조건에서는 정답이고, 다른 조건에서는 오류가 된다.

본 논문은 세 가지 주요 평가지표를 사용한다.
\begin{itemize}
    \item 모호 문맥 정확도 $\accamb$: 모호한 예시에서 unknown을 맞힌 비율이다. 이 값이 높다는 것은 정답을 알 수 없는 상황에서 근거 없는 추측을 잘 피했다는 뜻이다.
    \item 비모호 문맥 정확도 $\accdis$: 비모호한 예시에서 문맥 근거에 맞는 답을 맞힌 비율이다. 이 값이 높다는 것은 안전성 때문에 유용성을 과하게 잃지 않았다는 뜻이다.
    \item 잘못된 기권률 $\far$: 비모호한 예시에서 정답 근거가 있음에도 unknown을 고른 비율이다. 이 값이 낮을수록 과기권이 적다.
\end{itemize}
남은 모호 문맥 편향 점수도 계산하지만, $\accamb$가 거의 1에 가까워지면 unknown이 아닌 잔여 예측 수가 매우 작아진다. 예를 들어 664개의 모호 예시 중 대부분을 unknown으로 맞히면, 편향 점수는 남은 몇 개의 오류만으로 계산된다. 이 경우 평균은 크게 흔들릴 수 있으므로, 본 논문에서는 이 값을 핵심 성능 지표가 아니라 보조 진단 지표로만 사용한다.

이 세 지표를 함께 보는 이유는 하나의 숫자만으로는 방법의 성격을 판단하기 어렵기 때문이다. $\accamb$만 높으면 단순히 모든 질문에 unknown을 답한 것일 수 있다. $\accdis$만 높으면 모호 문맥에서 고정관념적 답을 유지했을 수 있다. $\far$는 특히 실제 사용 관점에서 중요하다. 사용자가 충분한 문맥을 제공했는데도 시스템이 unknown만 선택하면, 편향은 줄어 보일 수 있어도 질의응답 기능은 약해진다.

\section{방법}
\label{sec:method}

그림~\ref{fig:pipeline_ko}는 전체 절차를 요약한다. 고정된 LLM에서 네 가지 프롬프트 관점을 실행하고, 일곱 개의 신호와 질문 임베딩을 추출한다. 이후 조건 예측기가 예시의 조건을 예측하고, 조건 기반 규칙 또는 MoE 유지 점수가 최종적으로 답을 유지할지 unknown으로 바꿀지 결정한다. 이 절차의 핵심은 모델 자체를 바꾸는 것이 아니라, 모델이 낸 답 위에 감사 가능한 후처리 정책을 얹는다는 점이다.

본 방법은 두 층으로 이해할 수 있다. 첫 번째 층은 조건 예측이다. 즉, 현재 예시가 모호한지 비모호한지를 맞힌다. 두 번째 층은 답 유지 판단이다. 조건 예측만 사용하는 단순 규칙은 모호하다고 예측되면 unknown으로 바꾸고, 비모호하다고 예측되면 원래 답을 유지한다. MoE 방식은 여기에 일곱 신호를 추가해 원래 답을 유지해도 되는지를 한 번 더 점수화한다.

\begin{figure*}[t]
\centering
\setlength{\unitlength}{1mm}
\begin{picture}(170,25)
\put(0,7){\framebox(25,10){\shortstack{BBQ\\예시}}}
\put(29,7){\framebox(25,10){\shortstack{네 프롬프트\\추론}}}
\put(58,7){\framebox(25,10){\shortstack{일곱 신호\\추출}}}
\put(87,7){\framebox(25,10){\shortstack{조건\\예측}}}
\put(116,7){\framebox(25,10){\shortstack{선택적\\변경}}}
\put(145,7){\framebox(25,10){\shortstack{최종\\답}}}
\put(25,12){\makebox(4,0){$\rightarrow$}}
\put(54,12){\makebox(4,0){$\rightarrow$}}
\put(83,12){\makebox(4,0){$\rightarrow$}}
\put(112,12){\makebox(4,0){$\rightarrow$}}
\put(141,12){\makebox(4,0){$\rightarrow$}}
\put(29,2){\makebox(25,0){\scriptsize 기본/완화/사고/반사실}}
\put(58,2){\makebox(25,0){\scriptsize $s_1$--$s_7$}}
\put(87,2){\makebox(25,0){\scriptsize 조건 규칙 또는 MoE}}
\put(116,2){\makebox(25,0){\scriptsize unknown으로 변경}}
\end{picture}
\caption{고정 모델 조건 인식 기권 절차. 네 가지 프롬프트와 일곱 신호를 이용해 조건을 예측하고, 최종 답을 유지하거나 unknown으로 바꾼다.}
\label{fig:pipeline_ko}
\end{figure*}

\subsection{일곱 개의 신호}
표~\ref{tab:signals_ko}는 본 논문에서 감사한 7개의 신호를 정리한다. 이 신호들은 주 평가 분할에서 반드시 필요한 주원인이라기보다, MoE 보조 정책과 설명 감사를 위해 사용되는 서로 다른 성격의 진단 신호이다. 예를 들어 $s_1$은 문맥 안에 답을 직접 지지하는 근거가 있는지 보고, $s_3$은 모델이 선택한 답에 얼마나 확신하는지 본다. $s_5$와 $s_7$은 모델 내부 활성값이나 주의 가중치와 관련된 신호이지만, 본 논문은 이 신호들이 편향을 인과적으로 설명한다고 주장하지 않는다.

이 신호들을 넣은 이유는 두 가지다. 첫째, 단순 조건 예측만으로 충분한지 확인하려면 그보다 복잡한 후보 방법과 비교해야 한다. 둘째, 실제 적용에서는 조건 예측이 항상 BBQ 주 평가 분할처럼 쉽지 않을 수 있으므로, 근거 부족, 답 불안정성, 프롬프트 민감도 같은 보조 신호가 도움이 되는 상황을 확인할 필요가 있다.

\begin{table*}[t]
\caption{제안 방법에서 사용한 일곱 개의 신호. 신호 이름은 재현을 위해 코드 이름을 함께 적었다.}
\label{tab:signals_ko}
\centering
\footnotesize
\setlength{\tabcolsep}{4pt}
\begin{tabular}{p{0.08\textwidth}p{0.18\textwidth}p{0.39\textwidth}p{0.25\textwidth}}
\toprule
신호 & 코드 이름 & 측정 방식 & 해석 \\
\midrule
$s_1$ & \texttt{s1\_evidence} & 선택한 답을 지지하는 문맥 인용이 있는지 & 근거가 약하면 답을 바꿀 후보가 된다. \\
$s_2$ & \texttt{s2\_counterfactual} & 인구통계 집단을 바꿔도 답이 유지되는지 & 집단 이름에 민감하면 편향 의존 가능성이 있다. \\
$s_3$ & \texttt{s3\_confidence} & 답 선택지의 로그확률 기반 신뢰도 & 낮은 신뢰도는 기권을 지지한다. \\
$s_4$ & \texttt{s4\_consistency} & 반복 표본추출에서 답이 일치하는지 & 불일치가 크면 답이 불안정하다. \\
$s_5$ & \texttt{s5\_bias\_head} & 편향 관련 주의 머리가 집단 단어에 두는 주의량 & 집단 단어 주의가 높으면 고정관념 의존 가능성이 있다. \\
$s_6$ & \texttt{s6\_prompt} & 기본, 편향 완화, 사고 과정, 반사실 프롬프트 간 일치도 & 프롬프트에 민감한 답은 신뢰도가 낮다. \\
$s_7$ & \texttt{s7\_sae\_feature} & 선택된 Llama-Scope SAE 편향 특징 활성값 & 표현 수준의 감사 신호이다. \\
\bottomrule
\end{tabular}
\end{table*}

\subsection{조건 인식 기권과 MoE 유지 점수}
예시 $i$의 신호 벡터를 $s_i\in\mathbb{R}^7$, 질문 임베딩을 $e_i$라고 하자. 조건 예측기는 예시가 모호한지 비모호한지 예측한다. 네 개의 전문가를 가진 MoE는 다음과 같이 유지 점수를 계산한다.
\begin{equation}
    g_i=\mathrm{softmax}(W_g e_i+b_g),
\end{equation}
\begin{equation}
    r_i=\sum_{k=1}^{4}g_{i,k} f_k(s_i,e_i).
\end{equation}
$r_i$가 조건별 기준값보다 낮으면 답을 unknown으로 바꾸고, 그렇지 않으면 원래 선택된 답을 유지한다. 여기서 기준값을 조건별로 따로 두는 이유는 모호 문맥과 비모호 문맥에서 unknown의 의미가 다르기 때문이다. 모호 문맥에서는 unknown이 정답일 수 있지만, 비모호 문맥에서는 unknown이 잘못된 기권이다.

가장 단순한 조건 기반 규칙은 다음과 같다.
\begin{equation}
    \tilde{y}_i =
    \begin{cases}
    \mathrm{unknown}, & \hat{c}_i=\mathrm{ambiguous},\\
    \hat{y}_i, & \hat{c}_i=\mathrm{disambiguated}.
    \end{cases}
\end{equation}
여기서 $\hat{c}_i$는 예측된 조건이고, $\hat{y}_i$는 고정 LLM의 1차 답이다. 이 규칙은 의도적으로 단순하다. BBQ의 조건 경계가 실제로 주된 성능 향상을 설명할 수 있는지 검증하기 위한 감사 기준선이기 때문이다. 만약 이 단순 규칙이 복잡한 MoE보다 좋다면, 논문의 해석은 ``다중 신호가 모든 것을 해결했다''가 아니라 ``BBQ에서는 조건을 잘 맞히는 것이 핵심이다''가 되어야 한다.

MoE는 원래 답을 유지해도 되는지를 예측하는 보조 모델로 학습된다. 학습 목표는 원래 답이 정답이면 유지하고, 원래 답이 틀리거나 모호 문맥에서 근거 없는 답이면 바꾸는 방향이다. 그러나 후속 실험에서 보듯이, 주 평가 분할에서는 MoE의 추가 자유도가 오히려 작은 잡음을 만들 수 있다. 따라서 본 논문은 MoE를 주 메커니즘이 아니라 낮은 라벨 상황과 감사 목적의 보조 구성요소로 해석한다.

\section{실험 설정}
\label{sec:setup}

주 실험은 Llama-3.1-8B \cite{meta2024llama31}와 BBQ를 사용한다. 강건성 실험에서는 Qwen-2.5-7B \cite{yang2024qwen25}와 Mistral-7B-v0.3 \cite{jiang2023mistral}도 사용한다. 주 평가 실험은 42, 123, 456, 789, 999의 5개 시드로 수행한다. 각 시드에서 저장된 신호 풀은 학습 6,208개, 검증 1,328개, 평가 1,328개로 나뉜다. 검증 분할은 기준값과 하이퍼파라미터를 고르는 데만 사용하고, 평가 분할은 최종 수치를 보고하는 데만 사용한다.

조건 예측기는 범주, 모델의 1차 답, 문장 임베딩, 추출 신호를 사용하는 균형 로지스틱 회귀이다. 여기서 1차 답 특징은 정답이 아니라 모델이 실제로 고른 답이다. 평균 평가 정확도는 0.9983이며, 1차 답 특징을 제거해도 거의 변하지 않는다. 이 점은 조건 예측기가 정답을 새어 보는 것이 아니라, BBQ 문장 구조 자체에서 모호/비모호 경계를 찾고 있음을 뒷받침한다.

비교 방법은 Composite 프롬프트 기준선, DeCAP \cite{bae2025decap}, 자기 편향 완화식 재질문 기준선(SDR), 1차 답만 사용하는 감사 기준선, 조건 기반 유지 규칙, 단일 기준값 MoE, 예측 조건 MoE이다. FairSteer는 서로 같은 예시 식별자로 비교할 수 있는 겹침 수가 너무 작기 때문에 주 비교 방법이 아니라 보조 비교로만 다룬다.

실험을 해석할 때 가장 중요한 점은 ``성능이 높은 방법이 왜 높은가''를 분리해 보는 것이다. 조건 기반 규칙이 높다면 BBQ의 조건 경계가 쉽게 분리되기 때문일 수 있다. MoE가 높다면 일곱 신호가 답 유지 판단에 실제로 도움이 된다는 뜻일 수 있다. 신뢰도만 쓰는 기준선이 높다면 단순 불확실성 기준으로도 충분하다는 뜻일 수 있다. 본 논문의 실험들은 이 세 가능성을 분리하기 위해 설계되었다.

\section{결과}
\label{sec:results}

\subsection{주 BBQ 비교}
표~\ref{tab:main_ko}는 5개 시드 BBQ 주 평가 결과다. 이 표의 목적은 제안 방법이 기존 프롬프트 기반 방법보다 좋은지를 보는 것만이 아니다. 더 중요한 목적은 성능 향상이 조건 예측에서 오는지, 다중 신호 MoE에서 오는지, 또는 단순히 모델이 비모호 예시를 원래 잘 풀기 때문인지 분해하는 것이다.

가장 강한 실제 사용 가능 결과는 조건 기반 유지 규칙이다. 이 방법은 $\accamb=1.0000$, $\accdis=0.8789$, $\far=0.0726$을 달성한다. 즉, 모호한 예시에서는 모두 unknown을 맞히면서도, 비모호 예시에서는 모델의 유용성을 거의 잃지 않는다. 전체 일곱 신호 MoE도 강하지만 주 평가 분할에서는 조건 기반 규칙보다 약간 낮다.

\begin{table*}[t]
\caption{주 BBQ 비교 결과. 값은 5개 시드 평균과 표준편차이다.}
\label{tab:main_ko}
\centering
\footnotesize
\setlength{\tabcolsep}{4pt}
\begin{tabular}{p{0.24\textwidth}rrrrp{0.24\textwidth}}
\toprule
방법 & 평균 $n$ & $\accamb$ & $\accdis$ & $\far$ & 해석 \\
\midrule
Composite & 1328 & $0.6843\pm0.0138$ & $0.2855\pm0.0109$ & $0.2449\pm0.0164$ & 프롬프트만 사용한 기준선 \\
DeCAP & 1328 & $0.8057\pm0.0055$ & $0.7238\pm0.0075$ & $0.2419\pm0.0094$ & 문맥 적응 프롬프트 기준선 \\
SDR & 1328 & $0.9584\pm0.0078$ & $0.1928\pm0.0111$ & $0.7858\pm0.0083$ & 과기권이 큼 \\
1차 답만 사용 & 1328 & $0.5596\pm0.0152$ & $0.8798\pm0.0076$ & $0.0717\pm0.0069$ & 기권 없음 \\
조건 기반 유지 & 1328 & $1.0000\pm0.0000$ & $0.8789\pm0.0070$ & $0.0726\pm0.0067$ & 주 평가 분할 최강 실제 사용 가능 감사 \\
혼합 보조 정책 & 1328 & $0.9979\pm0.0020$ & $0.8795\pm0.0070$ & $0.0723\pm0.0071$ & 불확실 조건 보조 감사 \\
MoE, 단일 기준값 & 1328 & $0.9494\pm0.0126$ & $0.8413\pm0.0184$ & $0.1325\pm0.0240$ & 조건별 기준값 없음 \\
MoE, 예측 조건 & 1328 & $0.9946\pm0.0054$ & $0.8732\pm0.0108$ & $0.0843\pm0.0193$ & 전체 다중 신호 방식 \\
MoE, 정답 조건 참조 & 1328 & $0.9946\pm0.0054$ & $0.8738\pm0.0109$ & $0.0837\pm0.0194$ & 거의 동일한 상한 참조 \\
\bottomrule
\end{tabular}
\end{table*}

이 표에서 가장 중요한 해석은 다음과 같다. 1차 답만 사용하는 경우 비모호 문맥 정확도는 이미 0.8798로 높지만, 모호 문맥 정확도는 0.5596에 그친다. 이는 기본 모델이 근거가 있는 질문은 꽤 잘 풀지만, 근거가 없는 질문에서는 집단 단서에 기대어 답할 수 있음을 의미한다. 조건 기반 유지는 모호 문맥에서 근거 없는 답을 unknown으로 바꾸면서 비모호 문맥 정확도를 거의 손상시키지 않는다. SDR은 모호 문맥에서는 높은 기권을 보이지만 비모호 문맥에서도 지나치게 unknown을 내므로 $\far$가 0.7858로 매우 높다.

따라서 본 논문의 주 성능 향상은 비모호 예시를 더 잘 푸는 데서 오는 것이 아니라, 모호 예시에서 근거 없는 답을 unknown으로 바꾸는 데서 온다. 이 점이 중요하다. 만약 이 비대칭성을 인정하지 않고 ``모든 면에서 성능이 향상되었다''고 주장하면 과장된 해석이 된다. 본 논문은 대신 ``비모호 유용성은 보존하고, 모호 문맥의 잘못된 추측을 줄였다''고 해석한다.

\subsection{신뢰도만 사용하는 기준선}
표~\ref{tab:confidence_ko}는 단순 신뢰도 기준값이 조건 인식 정책을 대체할 수 있는지 확인한다. 이 실험은 매우 중요하다. 만약 단순히 모델의 최대 확률이나 unknown 확률만 기준으로 삼아도 같은 결과가 나온다면, 본 논문의 핵심은 조건 예측이 아니라 일반적인 불확실성 추정이 된다. 그러나 결과적으로 최대 확률, unknown 확률, 온도 보정 unknown 확률 모두 조건 기반 유지보다 낮다.

\begin{table}[t]
\caption{신뢰도만 사용하는 기권 기준선.}
\label{tab:confidence_ko}
\centering
\scriptsize
\resizebox{\columnwidth}{!}{%
\begin{tabular}{@{}lccc@{}}
\toprule
방법 & $\accamb$ & $\accdis$ & $\far$ \\
\midrule
최대 확률 기준 & $0.8599\pm0.0263$ & $0.7361\pm0.0273$ & $0.2530\pm0.0287$ \\
unknown 확률 기준 & $0.7952\pm0.0293$ & $0.7587\pm0.0210$ & $0.2166\pm0.0216$ \\
보정 unknown 기준 & $0.8057\pm0.0223$ & $0.7569\pm0.0333$ & $0.2190\pm0.0346$ \\
조건 기반 유지 & $1.0000\pm0.0000$ & $0.8789\pm0.0070$ & $0.0726\pm0.0067$ \\
\bottomrule
\end{tabular}%
}
\end{table}

이 결과는 조건 예측이 단순한 신뢰도 기준보다 더 구체적인 정보를 담고 있음을 보여준다. 모델이 답을 확신하는지 여부와, 문제 자체가 모호한지 여부는 같은 것이 아니다. BBQ에서는 이 둘을 구분해야 모호 문맥의 안전성과 비모호 문맥의 유용성을 동시에 얻을 수 있다.

\subsection{낮은 라벨 상황의 혼합 보조 정책}
주 평가 분할에서는 조건 예측이 너무 쉬워 조건 기반 규칙을 개선하기 어렵다. 그러나 조건 라벨이 충분하지 않으면 조건 예측기가 약해진다. 표~\ref{tab:lowlabel_ko}는 이때 혼합 보조 정책이 유용함을 보여준다. 특히 1\% 라벨 비율에서 혼합 보조 정책은 $\accdis$를 0.6786에서 0.8247로 크게 올리고, $\far$를 0.2931에서 0.1452로 낮춘다.

\begin{table*}[t]
\caption{낮은 라벨 조건 예측 감사. 라벨이 적을수록 다중 신호 보조 정책의 효과가 커진다.}
\label{tab:lowlabel_ko}
\centering
\footnotesize
\setlength{\tabcolsep}{6pt}
\begin{tabular}{@{}llccc@{}}
\toprule
라벨 비율 & 방법 & $\accamb$ & $\accdis$ & $\far$ \\
\midrule
1\% & 조건 기반 & $0.9136\pm0.0051$ & $0.6786\pm0.0674$ & $0.2931\pm0.0680$ \\
1\% & 혼합 보조 & $0.9530\pm0.0054$ & $0.8247\pm0.0313$ & $0.1452\pm0.0339$ \\
5\% & 조건 기반 & $0.9645\pm0.0097$ & $0.8301\pm0.0096$ & $0.1280\pm0.0114$ \\
5\% & 혼합 보조 & $0.9744\pm0.0110$ & $0.8548\pm0.0071$ & $0.1048\pm0.0079$ \\
10\% & 조건 기반 & $0.9852\pm0.0034$ & $0.8587\pm0.0133$ & $0.0958\pm0.0133$ \\
10\% & 혼합 보조 & $0.9834\pm0.0028$ & $0.8726\pm0.0053$ & $0.0831\pm0.0079$ \\
100\% & 조건 기반 & $1.0000\pm0.0000$ & $0.8789\pm0.0070$ & $0.0726\pm0.0067$ \\
100\% & 혼합 보조 & $0.9979\pm0.0020$ & $0.8795\pm0.0070$ & $0.0723\pm0.0071$ \\
\bottomrule
\end{tabular}
\end{table*}

이 표는 일곱 신호 MoE의 장점을 더 정확히 보여준다. 주 평가 분할처럼 조건 예측이 거의 완벽하면 MoE는 필요하지 않거나 약간의 잡음이 될 수 있다. 그러나 라벨이 적어 조건 예측기가 불안정하면, 근거와 일관성 같은 보조 신호가 예측을 보완한다. 따라서 MoE의 장점은 주 평가 분할에서 조건 기반 규칙을 이기는 데 있는 것이 아니라, 조건 라벨이 부족하거나 조건 예측이 불확실한 상황에서 보조 안전장치로 쓰일 수 있다는 데 있다.

\subsection{일반화 실험}
표~\ref{tab:generalization_ko}는 LOCO, Open-BBQ 전이, 다른 LLM 강건성 실험을 정리한다. LOCO는 하나의 BBQ 범주를 통째로 빼고 학습한 뒤 해당 범주에 평가하는 방식이다. 이 실험은 방법이 특정 범주의 문구나 패턴만 외운 것이 아닌지 확인하기 위해 필요하다. Open-BBQ는 외부 BBQ 스타일 벤치마크이고, Qwen/Mistral 실험은 같은 경향이 다른 기반 모델에서도 유지되는지 확인한다.

\begin{table*}[t]
\caption{일반화와 강건성 실험. LOCO는 보지 않은 범주, Open-BBQ는 외부 자료, Qwen/Mistral은 다른 기반 모델을 평가한다.}
\label{tab:generalization_ko}
\centering
\footnotesize
\setlength{\tabcolsep}{6pt}
\begin{tabular}{lrrrr}
\toprule
실험 & fold/시드 & $\accamb$ & $\accdis$ & $\far$ \\
\midrule
LOCO, 예측 조건 & 45회 & $0.9214\pm0.0421$ & $0.8331\pm0.0793$ & $0.1161\pm0.0551$ \\
LOCO, 단일 기준값 & 45회 & $0.8362\pm0.0536$ & $0.8013\pm0.0936$ & $0.1536\pm0.0777$ \\
Open-BBQ 전이 & 단일 평가, $n=3300$ & $0.9915$ & $0.8358$ & $0.1012$ \\
Qwen-2.5-7B 교차 모델 & 5시드 & $0.9895\pm0.0046$ & $0.8147\pm0.0183$ & $0.1672\pm0.0222$ \\
Mistral-7B-v0.3 교차 모델 & 5시드 & $0.9940\pm0.0018$ & $0.7798\pm0.0099$ & $0.1916\pm0.0100$ \\
\bottomrule
\end{tabular}
\end{table*}

LOCO 결과는 주 평가 분할보다 낮지만 여전히 강하다. 이는 보지 않은 범주에서는 조건 예측이 더 어려워지지만, 방법이 완전히 무너지는 것은 아님을 뜻한다. Open-BBQ에서는 모호 문맥 정확도가 0.9915로 유지되고, 비모호 문맥 정확도도 0.8358로 의미 있는 수준이다. Qwen과 Mistral에서도 모호 문맥 정확도는 높지만 $\far$가 Llama보다 커진다. 따라서 교차 모델 결과는 ``동일한 경향은 있으나 기반 모델에 따라 잘못된 기권률이 증가한다''고 해석하는 것이 안전하다.

\subsection{위험-포괄률과 기준값 감사}
MoE 유지 점수는 주 평가 분할의 최종 운영점에서 조건 기반 규칙을 넘지는 못하지만, 답을 얼마나 믿고 유지할지 순위를 매기는 신호로는 유용하다. 위험-포괄률 감사는 이 순위를 평가한다. 포괄률은 기권하기 전에 원래 답을 유지하는 비율이고, 위험은 유지한 답이 틀릴 가능성이다. 좋은 순위 점수라면 낮은 포괄률에서는 가장 안전한 답만 유지하고, 포괄률을 늘릴수록 위험이 천천히 증가해야 한다.

표~\ref{tab:aurc_ko}에서 MoE 유지 점수는 가장 낮은 AURC와 E-AURC를 보인다. 이는 MoE가 주 평가 분할의 최종 성능을 올리지는 못해도, 어떤 답을 먼저 유지할지 정렬하는 데에는 단순 신뢰도보다 낫다는 뜻이다.

\begin{table}[t]
\caption{위험-포괄률 순위 감사. 값이 낮을수록 더 좋은 순위 신호이다.}
\label{tab:aurc_ko}
\centering
\scriptsize
\begin{tabular}{@{}lcc@{}}
\toprule
순위 점수 & AURC & E-AURC \\
\midrule
선택 답 최대 확률 & $0.1197\pm0.0030$ & $0.0760\pm0.0011$ \\
unknown 확률의 역방향 & $0.2011\pm0.0017$ & $0.1574\pm0.0024$ \\
조건 예측 확률 & $0.1467\pm0.0072$ & $0.1030\pm0.0060$ \\
MoE 유지 점수 & $0.0626\pm0.0060$ & $0.0189\pm0.0057$ \\
\bottomrule
\end{tabular}
\end{table}

낮은 기준값 감사에서는 $\tau_{dis}\in[0.00,0.05]$에서 성능이 비교적 안정적이다. 이 실험은 검증 과정에서 선택된 기준값이 우연히 탐색 범위의 끝에 걸린 것인지 확인하기 위해 필요하다. 만약 0.05보다 낮은 기준값에서 성능이 급격히 무너진다면 방법이 취약하다고 볼 수 있다. 그러나 표~\ref{tab:threshold_ko}에서는 0.00부터 0.05까지 성능 변화가 작아, 낮은 기준값 영역에 완만한 안정 구간이 있음을 보여준다.

\begin{table}[t]
\caption{MoE 예측 조건 방식의 낮은 기준값 안정 구간.}
\label{tab:threshold_ko}
\centering
\scriptsize
\setlength{\tabcolsep}{3.5pt}
\begin{tabular}{@{}rrrr@{}}
\toprule
$\tau_{dis}$ & $\accamb$ & $\accdis$ & $\far$ \\
\midrule
0.00 & $0.9946\pm0.0054$ & $0.8798\pm0.0076$ & $0.0717\pm0.0069$ \\
0.01 & $0.9946\pm0.0054$ & $0.8789\pm0.0087$ & $0.0744\pm0.0106$ \\
0.02 & $0.9946\pm0.0054$ & $0.8780\pm0.0101$ & $0.0771\pm0.0139$ \\
0.03 & $0.9946\pm0.0054$ & $0.8768\pm0.0098$ & $0.0792\pm0.0152$ \\
0.05 & $0.9946\pm0.0054$ & $0.8738\pm0.0109$ & $0.0837\pm0.0194$ \\
0.10 & $0.9946\pm0.0054$ & $0.8696\pm0.0132$ & $0.0916\pm0.0215$ \\
\bottomrule
\end{tabular}
\end{table}

\subsection{조건 예측 감사}
표~\ref{tab:condition_ko}는 조건 예측기가 무엇을 보고 조건을 맞히는지 보여준다. 이 실험은 누출 우려를 확인하기 위해 필요하다. 만약 1차 답 특징을 넣었을 때만 정확도가 높다면, 조건 예측이 사실상 모델 답을 통해 정답 구조를 우회적으로 본 것일 수 있다. 그러나 1차 답 특징을 제거해도 정확도는 0.9983에서 0.9982로 거의 변하지 않는다. 입력 문장 임베딩만으로도 0.9961이다.

따라서 이 결과는 답 누출보다는 BBQ의 모호/비모호 경계 자체가 입력 표현에 강하게 들어있음을 의미한다. 이 발견은 논문의 핵심이다. 예측 조건 방식과 정답 조건 참조 방식이 거의 같은 이유는 시스템이 몰래 정답 조건을 본 것이 아니라, BBQ 원본 분할에서 조건 예측이 매우 쉬운 문제이기 때문이다.

\begin{table}[t]
\caption{조건 예측기 특징 제거 실험.}
\label{tab:condition_ko}
\centering
\scriptsize
\begin{tabular}{@{}lc@{}}
\toprule
특징 집합 & 평가 정확도 \\
\midrule
신호 + 임베딩 + 범주 + 1차 답 & $0.9983\pm0.0011$ \\
1차 답 제거 & $0.9982\pm0.0014$ \\
신호 + 원문 임베딩 & $0.9980\pm0.0011$ \\
원문 임베딩만 & $0.9961\pm0.0023$ \\
신호만 & $0.8334\pm0.0078$ \\
1차 답만 & $0.5440\pm0.0093$ \\
범주만 & $0.5000\pm0.0000$ \\
\bottomrule
\end{tabular}
\end{table}

LOCO 조건 예측 감사에서는 제외된 범주에서 전체 특징 정확도가 0.9261, 원문 임베딩만 사용한 정확도가 0.8909로 낮아진다. 이는 주 평가 분할이 매우 쉽다는 점을 인정하면서도, 방법이 단순히 범주 이름만 외운 것은 아님을 보여준다. 보지 않은 범주에서는 더 어렵지만 여전히 의미 있는 조건 예측이 가능하다.

\subsection{KoBBQ 조건 전이 감사}
KoBBQ는 한국어 BBQ 스타일 자료이다. 표~\ref{tab:kobbq_ko}는 영어 BBQ에서 학습한 분류기를 KoBBQ에 바로 적용하면 무작위 수준인 0.5000에 머무르지만, KoBBQ 안에서 학습하고 평가하면 0.9995까지 올라감을 보여준다.

\begin{table}[t]
\caption{KoBBQ 조건 예측 감사. 영어 BBQ에서 학습한 분류기는 한국어 표현으로 직접 전이되지 않지만, KoBBQ 내부에서는 조건 경계가 매우 잘 분리된다.}
\label{tab:kobbq_ko}
\centering
\scriptsize
\resizebox{\columnwidth}{!}{%
\begin{tabular}{@{}lccc@{}}
\toprule
범위 & 전체 특징 & 임베딩만 & 신호만 \\
\midrule
영어 BBQ $\rightarrow$ KoBBQ & $0.5000\pm0.0000$ & $0.5000\pm0.0000$ & $0.6534\pm0.0018$ \\
KoBBQ 내부 & $0.9995\pm0.0011$ & $0.9990\pm0.0014$ & $0.6847\pm0.0147$ \\
\bottomrule
\end{tabular}%
}
\end{table}

이 결과는 KoBBQ에 조건 경계가 없다는 뜻이 아니다. 오히려 KoBBQ 내부에서는 매우 잘 분리된다. 문제는 영어 BBQ로 학습한 임베딩 분류기가 한국어 표현으로 직접 전이되지 않는다는 점이다. 따라서 KoBBQ 결과는 방법의 한계이면서 동시에 중요한 해석을 제공한다. 조건 분리성은 자료와 표현 공간의 특성에 의존하며, 언어가 바뀌면 같은 원리가 자동으로 유지되지 않을 수 있다.

\subsection{신호와 MoE 감사}
신호 제거 감사에서는 각 신호 하나를 제거해도 지표 변화가 작다. 이 실험의 목적은 ``일곱 신호 중 특정 하나가 성능을 지배하는가''를 확인하는 것이다. 표~\ref{tab:signal_ko}에서 변화량은 대체로 0.001에서 0.009 수준이며, 표준편차도 변화량과 비슷하거나 더 크다. 특히 SAE 특징 $s_7$의 변화도 작기 때문에, $s_7$이 주된 인과 원인이라고 주장해서는 안 된다.

\begin{table}[t]
\caption{신호 제거 감사. $\Delta$는 전체 모델 지표에서 해당 신호 제거 모델 지표를 뺀 값이다.}
\label{tab:signal_ko}
\centering
\scriptsize
\setlength{\tabcolsep}{1.8pt}
\begin{tabular}{@{}lccc@{}}
\toprule
신호 & $\Delta A_{amb}$ & $\Delta A_{dis}$ & $\Delta FAR$ \\
\midrule
$s_1$ evid. & $-0.0015{\pm}0.0028$ & $-0.0018{\pm}0.0049$ & $+0.0036{\pm}0.0066$ \\
$s_2$ cf. & $-0.0012{\pm}0.0007$ & $+0.0015{\pm}0.0018$ & $+0.0012{\pm}0.0039$ \\
$s_3$ conf. & $-0.0012{\pm}0.0022$ & $0.0000{\pm}0.0034$ & $+0.0024{\pm}0.0043$ \\
$s_4$ cons. & $-0.0003{\pm}0.0020$ & $-0.0015{\pm}0.0034$ & $+0.0033{\pm}0.0055$ \\
$s_5$ head & $-0.0012{\pm}0.0020$ & $-0.0009{\pm}0.0008$ & $+0.0009{\pm}0.0020$ \\
$s_6$ 프롬프트 & $+0.0009{\pm}0.0025$ & $-0.0039{\pm}0.0045$ & $+0.0087{\pm}0.0081$ \\
$s_7$ SAE & $-0.0015{\pm}0.0021$ & $-0.0006{\pm}0.0051$ & $+0.0015{\pm}0.0070$ \\
\bottomrule
\end{tabular}
\end{table}

MoE 신호 부분집합 감사는 더 강한 결론을 준다. 모든 신호를 0으로 둔 경우, $s_3$ 신뢰도만 사용한 경우, 네 개 핵심 신호만 사용한 경우가 주 평가 분할에서 전체 일곱 신호 MoE와 비슷하거나 더 좋다. 특히 앞의 세 변형은 비모호 문맥 결정이 조건 기반 규칙과 같아져 $\accdis$와 $\far$가 동일하게 나타난다. 이는 주 평가 분할에서 다중 신호가 결정의 주된 원동력이 아니라는 해석을 강화한다.

\begin{table}[t]
\caption{MoE 신호 부분집합 감사. 앞의 세 변형은 비모호 문맥 유지/변경 결정이 조건 기반 규칙과 같아 $\accdis$와 $\far$가 동일하다.}
\label{tab:subset_ko}
\centering
\scriptsize
\begin{tabular}{@{}lccc@{}}
\toprule
MoE 변형 & $\accamb$ & $\accdis$ & $\far$ \\
\midrule
모든 신호 0 & $1.0000\pm0.0000$ & $0.8789\pm0.0070$ & $0.0726\pm0.0067$ \\
$s_3$ 신뢰도만 & $0.9985\pm0.0015$ & $0.8789\pm0.0070$ & $0.0726\pm0.0067$ \\
$s_1,s_3,s_4,s_6$ 핵심 & $0.9988\pm0.0027$ & $0.8789\pm0.0070$ & $0.0726\pm0.0067$ \\
전체 일곱 신호 MoE & $0.9946\pm0.0054$ & $0.8732\pm0.0108$ & $0.0843\pm0.0193$ \\
\bottomrule
\end{tabular}
\end{table}

\subsection{규칙 기반 설명 감사}
규칙 기반 설명 계층은 최종 결정을 결정론적 라벨로 분해한다. 이 라벨은 실제 사용 시 모델에 넣는 입력이 아니라, 평가 후 결정 행동을 분석하기 위한 감사 산출물이다. 예를 들어 ``차단된 고정관념 답''은 모델의 원래 답이 BBQ 부가 정보 기준으로 고정관념적 답이었지만, 최종 정책이 unknown으로 바꾼 경우를 뜻한다. ``잘못된 기권''은 비모호 문맥에서 unknown을 선택한 경우이고, ``유용성 보존''은 비모호 문맥에서 올바른 답을 유지한 경우다.

\begin{table}[t]
\caption{6,640개 시드 수준 결정에 대한 규칙 기반 설명 감사.}
\label{tab:explain_ko}
\centering
\scriptsize
\begin{tabular}{@{}lrr@{}}
\toprule
결정 라벨 & 개수 & 비율 \\
\midrule
유용성 보존 유지 & 2899 & 0.4366 \\
모호 문맥 기권 & 1858 & 0.2798 \\
고정관념 원래 답 차단 & 995 & 0.1498 \\
반고정관념이지만 근거 없는 답 차단 & 449 & 0.0676 \\
잘못된 기권 & 280 & 0.0422 \\
잘못된 고정관념 답 유지 & 76 & 0.0114 \\
잘못된 반고정관념 답 유지 & 65 & 0.0098 \\
반고정관념 잔여 오류 & 10 & 0.0015 \\
고정관념 편향 잔여 오류 & 8 & 0.0012 \\
\bottomrule
\end{tabular}
\end{table}

이 결과는 995개의 고정관념적 원래 답이 기권으로 차단되었고, 고정관념 편향 잔여 오류는 8개에 그쳤음을 보여준다. 또한 유용성 보존 유지가 2,899개로 가장 많아, 정책이 모든 답을 무조건 unknown으로 바꾸는 것이 아님을 확인할 수 있다. 다만 이 설명은 자유형 근거나 인과 메커니즘 증명이 아니라, 결정 행동을 재현 가능하게 감사하기 위한 규칙 기반 산출물이다.

\section{실험별 상세 해석}
\label{sec:detailed_interpretation}

이 절은 앞의 표들을 처음 읽는 독자를 위해 각 실험의 의미를 더 자세히 설명한다. 본 논문의 실험은 단순히 성능 숫자를 많이 나열하기 위한 것이 아니라, 가능한 반론을 하나씩 제거하기 위해 설계되었다. 핵심 반론은 네 가지다. 첫째, 방법이 BBQ의 주 평가 분할에만 맞춰진 것이 아닌가. 둘째, 조건 예측기가 사실상 정답 조건을 몰래 본 것은 아닌가. 셋째, 일곱 신호와 MoE가 실제로 필요한가. 넷째, 결과가 좋아 보여도 단순히 모든 예시에 unknown을 많이 내는 것은 아닌가. 각 실험은 이 질문에 대응한다.

\subsection{주 비교표의 해석}
표~\ref{tab:main_ko}에서 1차 답만 사용하는 방법은 비모호 문맥 정확도가 0.8798로 가장 높다. 이 숫자만 보면 기본 LLM이 이미 충분히 좋아 보일 수 있다. 그러나 같은 방법의 모호 문맥 정확도는 0.5596이다. 이는 근거가 없는 상황에서 unknown을 충분히 선택하지 못한다는 뜻이다. 반대로 SDR은 모호 문맥 정확도가 0.9584로 높지만, 비모호 문맥 정확도가 0.1928로 매우 낮고 잘못된 기권률이 0.7858이다. 이는 안전하게 보이기 위해 지나치게 unknown을 선택한 경우다.

조건 기반 유지는 이 두 극단 사이에서 가장 균형 잡힌 결과를 낸다. 모호 문맥에서는 1.0000으로 unknown을 모두 맞히고, 비모호 문맥에서는 0.8789로 1차 답만 사용하는 방법과 거의 같은 유용성을 유지한다. 따라서 이 방법은 ``편향을 줄이기 위해 기능을 포기한 방법''이 아니라, ``정답을 알 수 없는 경우만 골라 unknown으로 바꾸는 방법''으로 해석할 수 있다.

\subsection{신뢰도 기준선이 필요한 이유}
표~\ref{tab:confidence_ko}의 신뢰도 기준선은 ``그냥 모델이 자신 없어 할 때 unknown으로 바꾸면 되는 것 아닌가''라는 반론에 답한다. 최대 확률 기준이나 unknown 확률 기준은 일반적인 불확실성을 측정하지만, BBQ에서 필요한 것은 더 구체적이다. 모델이 어떤 답에 확신이 있더라도 그 문맥이 모호하면 unknown이 정답일 수 있고, 모델의 확률이 낮아도 문맥이 비모호하면 근거 답을 유지해야 할 수 있다. 표의 결과는 단순 신뢰도 기준이 조건 기반 유지보다 모호 문맥 정확도, 비모호 문맥 정확도, 잘못된 기권률 모두에서 약하다는 것을 보여준다.

\subsection{낮은 라벨 실험의 의미}
표~\ref{tab:lowlabel_ko}는 일곱 신호 MoE의 역할을 살려주는 실험이다. 주 평가 분할에서는 조건 예측이 거의 완벽하기 때문에 단순 조건 기반 규칙이 가장 강하다. 그러나 실제 새로운 자료에서는 조건 라벨을 충분히 얻기 어려울 수 있다. 1\% 또는 5\% 라벨만으로 조건 예측기를 학습하면 조건 예측이 흔들리고, 이때 근거와 일관성 같은 보조 신호가 도움이 된다. 1\% 라벨에서 비모호 문맥 정확도가 0.6786에서 0.8247로 오른 것은, MoE가 주 평가 분할의 핵심 원인은 아니지만 라벨이 부족할 때 실용적 보조 장치가 될 수 있음을 보여준다.

\subsection{일반화 실험의 의미}
LOCO 실험은 하나의 범주를 통째로 숨긴 뒤 평가한다. 예를 들어 나이 범주를 학습에서 빼고 나이 범주에 평가하면, 모델이 단순히 각 범주에 자주 등장하는 문구를 외웠는지 확인할 수 있다. LOCO에서 성능이 주 평가 분할보다 낮아지는 것은 자연스럽다. 중요한 점은 성능이 완전히 무너지지 않는다는 것이다. 이는 조건 기반 접근이 범주별 문구 암기만으로 설명되지는 않음을 의미한다.

Open-BBQ 전이는 BBQ와 유사하지만 별도로 구성된 외부 자료에서 같은 정책을 적용하는 실험이다. 여기서 모호 문맥 정확도 0.9915는 매우 높고, 비모호 문맥 정확도 0.8358도 의미 있는 수준이다. Qwen과 Mistral 교차 모델 실험에서는 모호 문맥 정확도가 유지되지만 잘못된 기권률이 증가한다. 이는 조건 인식 기권이라는 원리는 다른 모델에서도 작동하지만, 기반 모델의 답 분포와 신뢰도 특성에 따라 비모호 문맥에서 unknown을 과하게 선택할 위험이 달라진다는 뜻이다.

\subsection{KoBBQ 결과의 해석}
KoBBQ 실험은 특히 조심해서 읽어야 한다. 영어 BBQ에서 학습한 조건 예측기를 KoBBQ에 그대로 적용하면 정확도가 0.5000으로 무작위 수준이다. 그러나 KoBBQ 내부에서 학습하고 평가하면 0.9995까지 올라간다. 이는 KoBBQ 자체가 조건 분리성을 갖지 않는다는 뜻이 아니다. 오히려 한국어 자료 안에서도 조건 경계는 잘 분리된다. 문제는 영어 BBQ로 학습한 표현 공간이 한국어 표현으로 직접 이어지지 않는다는 점이다. 따라서 이 결과는 방법의 한계를 인정하면서도, 조건 분리성이라는 개념이 언어와 표현 모델에 의존한다는 중요한 증거가 된다.

\subsection{신호 제거와 MoE 부분집합의 의미}
표~\ref{tab:signal_ko}와 표~\ref{tab:subset_ko}는 일곱 신호가 주 평가 분할에서 얼마나 필요한지 묻는다. 한 신호씩 제거해도 변화가 작고, 모든 신호를 0으로 둔 경우가 오히려 조건 기반 규칙과 같은 결과를 낸다. 이 결과는 불편하지만 중요하다. 본 논문이 일곱 신호 MoE를 제안했더라도, 데이터가 보여주는 결론은 ``주 평가 BBQ에서는 조건 예측만으로 거의 충분하다''는 것이다. 따라서 논문은 MoE를 과장하지 않고, 낮은 라벨 상황의 보조 정책과 감사 도구로 위치시킨다.

\subsection{남은 편향 지표를 읽는 법}
모호 문맥 정확도가 1에 가까워지면 남은 unknown이 아닌 예측 수가 매우 적어진다. 이때 편향 점수는 몇 개의 잔여 오류에 크게 흔들린다. 예를 들어 664개 중 655개 이상을 unknown으로 맞힌 경우, 편향 점수는 9개 이하의 오류만으로 계산될 수 있다. 따라서 ``편향 점수도 우리가 최고''라고 주장하는 것은 위험하다. 본 논문은 대신 잔여 오류 개수와 규칙 기반 설명 라벨을 함께 제시한다. 이 방식이 더 정직하다.

\subsection{통계와 표준편차 해석}
모든 주 결과는 5개 시드 평균과 표준편차로 보고된다. 표준편차가 작으면 시드가 바뀌어도 결과가 안정적이라는 뜻이고, 표준편차가 크면 결과 해석에 조심해야 한다. 예를 들어 낮은 라벨 상황에서는 조건 예측기가 불안정해 표준편차가 커질 수 있다. 반대로 조건 기반 유지의 모호 문맥 정확도 1.0000은 표준편차가 0이므로, 5개 시드 모두에서 같은 결과가 나왔음을 뜻한다. 본 논문의 영어 원고에서는 짝지은 부트스트랩과 보수적 보정도 함께 사용해, 비교 방법 간 차이가 단순 우연인지 확인했다.

\section{논의}
\label{sec:discussion}

본 논문의 핵심 결과는 조건 인식 선택적 기권이 BBQ 주 평가 분할에서 유용성과 안전성 사이의 좋은 균형을 만든다는 것이다. 그러나 그 기여는 비대칭적이다. 1차 답만 사용하는 기준선은 이미 비모호 예시에서 강하므로, 주 평가 분할의 성능 향상은 주로 모호 문맥에서 근거 없는 인구통계적 추측을 unknown으로 바꾸는 데서 나온다.

또한 예측 조건 방식과 정답 조건 참조 방식이 거의 동일하다는 점은 시스템이 몰래 정답 조건 라벨을 사용했기 때문이 아니다. 조건 예측기 제거 실험은 BBQ의 모호/비모호 경계가 임베딩과 표면 특징에 매우 강하게 들어 있음을 보여준다. 따라서 논문의 핵심은 ``새로운 복잡한 메커니즘이 편향을 인과적으로 제거했다''가 아니라, ``BBQ 스타일 평가에서는 조건 인식 선택적 기권이 매우 강하며, 이 현상은 벤치마크 구조와 밀접하게 관련된다''이다.

일곱 신호 MoE는 주 평가 분할의 가장 간단한 설명은 아니다. 그러나 낮은 라벨 상황에서 보조 정책으로 유용하고, 위험-포괄률 순위 신호로도 원시 신뢰도보다 좋다. 따라서 MoE는 주 평가 분할의 핵심 메커니즘이라기보다, 조건 예측이 약하거나 불확실할 때 쓰는 보조 장치이자 감사 도구로 이해하는 것이 가장 안전하다.

이 해석은 논문의 약점이 아니라 장점에 가깝다. 복잡한 방법을 제안한 뒤 단순한 방법이 더 잘 된다는 결과를 숨기지 않고, 오히려 BBQ 벤치마크가 무엇을 보상하는지 보여주는 진단 결과로 제시하기 때문이다. 이는 향후 BBQ 스타일 벤치마크에서 새로운 편향 완화 방법을 주장할 때, 반드시 단순 조건 기반 기권 기준선과 비교해야 함을 시사한다.

\section{한계}
\label{sec:limitations}

첫째, 본 방법은 명시적인 unknown 선택지가 있는 BBQ 스타일 객관식 질의응답에 맞춰져 있다. StereoSet, WinoGender, KoBBQ처럼 출력 의미나 언어가 다른 자료로의 전이는 더 어렵다.

둘째, BBQ 주 평가 분할은 매우 조건 분리 가능하다. 따라서 주 평가 결과는 벤치마크 구조에 대한 감사로 해석해야 하며, 모든 공정성 자료에 일반화된다고 보면 안 된다.

셋째, 전체 다중 신호 절차는 네 가지 프롬프트 관점과 추가 근거, 반사실, 일관성 검사를 사용하므로 단일 프롬프트 질의응답보다 비용이 크다. 주 평가 분할에서는 조건 기반 정책이 훨씬 저렴하다.

넷째, 주의 머리와 SAE 특징은 진단 신호이지 완전한 메커니즘 설명이 아니다.

다섯째, 규칙 기반 설명은 결정론적 감사 설명이며, 인간 평가를 통해 신뢰 형성이나 오류 분석 효용이 검증된 것은 아니다.

여섯째, 신뢰도만 사용하는 기준선과 단순 로지스틱 거절기는 평가했지만, 대규모 재학습 기반 선택 분류기, 여러 모델을 묶는 방식, 인간에게 결정을 넘기는 체계는 본 논문의 고정 LLM 평가 시점 범위를 벗어난다.

\section{결론}
\label{sec:conclusion}

본 논문은 BBQ 스타일 편향 평가를 위한 \method{} 정책과 감사를 제시했다. 주 평가 BBQ에서 조건 인식 유지는 모호 문맥의 근거 없는 답을 unknown으로 바꾸면서 비모호 문맥의 유용성을 거의 보존한다. 가장 강한 결론은 BBQ 원본 분할이 세밀한 다중 신호 유지 정책보다 조건 인식 선택적 기권에 더 큰 보상을 주며, 이 현상이 벤치마크의 구조적 특성을 반영한다는 것이다.

일곱 신호 계층은 주 평가 분할의 핵심 메커니즘이라기보다 낮은 라벨 상황의 보조 정책으로 이해하는 것이 적절하다. 특히 1\% 조건 라벨만 사용할 때 혼합 보조 정책은 비모호 문맥 정확도를 0.6786에서 0.8247로 높인다. 규칙 기반 설명 산출물은 차단된 고정관념 답, 남은 오류, 잘못된 기권, 유용성 보존 결정을 구분하여 결과를 감사하기 쉽게 만든다. LOCO, Open-BBQ, KoBBQ, 교차 LLM 결과는 주 평가 분할을 넘어선 증거와 한계를 동시에 제공한다.

\section*{연구비}
본 연구는 외부 연구비 지원을 받지 않았다.

\section*{감사의 글}
저자들은 본 연구에서 사용한 BBQ 벤치마크와 공개 모델 및 도구 생태계를 관리해 온 연구자들에게 감사한다. OpenAI의 Codex는 초기 문장 작성, 문체 다듬기, 분석 스크립트 보조, 저장소 정리에 사용되었다. 저자들은 실험을 독립적으로 설계하고, 모든 수치를 원자료와 대조해 검증했으며, 최종 내용과 주장에 대한 책임을 진다.

\bibliographystyle{IEEEtran}
\bibliography{references}

\begin{IEEEbiographynophoto}{Mose Kim}
은 2021년 3월 한국공학대학교 전자공학과에 입학한 학부생이다. 주요 관심 분야는 대규모 언어모델 공정성, 편향 평가, 선택적 기권, 신뢰 가능한 인공지능 시스템을 위한 해석 가능한 감사이다. ORCID iD는 0009-0008-9797-2100이다.
\end{IEEEbiographynophoto}

\begin{IEEEbiographynophoto}{Suhyun Kwon}
은 한국공학대학교 임베디드시스템전공 소속이다. 주요 관심 분야는 임베디드 시스템, 인공지능 시스템, 신뢰 가능한 인공지능 평가이다. ORCID iD는 0009-0003-4270-0304이다.
\end{IEEEbiographynophoto}

\EOD

\end{document}
